% GENERATED FILE. Edit canonical lessons, then run npm run generate:books.
% canonical-source-hash: fnv1a64:3b1005d3acc6b511
% canonical-lessons: FA-C436-eskele, FA-C436-nakhoda, FA-C436-arshe, FA-C436-lenj, FA-C436-langar-andakhtan

\chapter{By boat}
\label{ch:fa-by-boat}
\hlchaptermodality{436}

\begin{chapteropening}
\textbf{By the end of this chapter:} \emph{I can talk about travelling by sea: the pier, the captain, the deck, a dhow and dropping anchor.}

Complete the last lesson of chapter 436: I can talk about travelling by sea: the pier, the captain, the deck, a dhow and dropping anchor.
\end{chapteropening}

\section[\texorpdfstring{\fa{اسکله} (eskele)}{اسکله (eskele)}]{\fa{اسکله} (eskele) — a pier, a jetty, a quay}

\label{lesson:FA-C436-eskele}



\begin{quote}
Before the new one: say the Persian for a porter, then the Persian for to give someone a lift.
\end{quote}

\subsection*{You'll want to know: \fa{اسکله}}
\textbf{\fa{اسکله}} (\emph{eskele}) — \textquotedblleft{}a pier, a jetty, a quay\textquotedblright{}.

\begin{quote}
\fa{صبح} \fa{قایق} \fa{از} \fa{اسکله} \fa{حرکت} \fa{کرد}.
\end{quote}

\begin{quote}
\emph{Sobh qâyeq az eskele harekat kard.}
\end{quote}

\begin{quote}
In the morning the boat left the pier.
\end{quote}

Ports on the Persian Gulf have long \emph{eskele} for boats to the islands.

\begin{cousinweb}
\emph{eskele} came into Persian from Turkish \emph{iskele} 'a landing place'; EtymDB lists the Turkish word without a gloss. The Turkish word is usually traced to Italian \emph{scala} 'ladder, landing'.
\end{cousinweb}

\begin{grammarlens}[title={What you've built}]
One more word for telling, explaining and describing at B1.
\end{grammarlens}

\subsection*{Your turn}
\begin{itemize}
\raggedright
  \item \emph{Say it:} \emph{eskele}
  \item \emph{Say it:} \emph{eskele}, once more
  \item \emph{Say it:} a sentence of your own with \emph{eskele}
  \item \emph{Recall:} say the Persian for a porter, then the Persian for to give someone a lift, then say \emph{eskele} again
\end{itemize}

\begin{tcolorbox}[breakable,colback=teal!4,colframe=teal!35!black,title={Before you move on}]
What does \fa{اسکله} mean? (\textquotedblleft{}A pier, a jetty, a quay\textquotedblright{}.) Say it once more.
\end{tcolorbox}
\section[\texorpdfstring{\fa{ناخدا} (nâkhodâ)}{ناخدا (nâkhodâ)}]{\fa{ناخدا} (nâkhodâ) — a captain (of a ship), a skipper}

\label{lesson:FA-C436-nakhoda}



\begin{quote}
Before the new one: say the Persian for to give someone a lift, then the Persian for a pier.
\end{quote}

\subsection*{You'll want to know: \fa{ناخدا}}
\textbf{\fa{ناخدا}} (\emph{nâkhodâ}) — \textquotedblleft{}a captain (of a ship), a skipper\textquotedblright{}.

\begin{quote}
\fa{ناخدا} \fa{سی} \fa{سال} \fa{روی} \fa{دریا} \fa{کار} \fa{کرده} \fa{است}.
\end{quote}

\begin{quote}
\emph{Nâkhodâ si sâl ru-ye daryâ kâr karde ast.}
\end{quote}

\begin{quote}
The captain has worked at sea for thirty years.
\end{quote}

Literally 'ship-master'. The navy rank is also \emph{nâkhodâ}.

\begin{cousinweb}
\emph{nâkhodâ} is a native compound, which EtymDB analyses as \emph{nâv} 'ship' and \emph{khodâ} 'chief, master': the master of the ship. \emph{Khodâ} is usually traced to Middle Persian \emph{xwadāy} 'lord'.
\end{cousinweb}

\begin{grammarlens}[title={What you've built}]
One more word for telling, explaining and describing at B1.
\end{grammarlens}

\subsection*{Your turn}
\begin{itemize}
\raggedright
  \item \emph{Say it:} \emph{nâkhodâ}
  \item \emph{Say it:} \emph{nâkhodâ}, once more
  \item \emph{Say it:} a sentence of your own with \emph{nâkhodâ}
  \item \emph{Recall:} say the Persian for to give someone a lift, then the Persian for a pier, then say \emph{nâkhodâ} again
\end{itemize}

\begin{tcolorbox}[breakable,colback=teal!4,colframe=teal!35!black,title={Before you move on}]
What does \fa{ناخدا} mean? (\textquotedblleft{}A captain (of a ship), a skipper\textquotedblright{}.) Say it once more.
\end{tcolorbox}
\section[\texorpdfstring{\fa{عرشه} ('arshe)}{عرشه ('arshe)}]{\fa{عرشه} ('arshe) — a deck (of a ship)}

\label{lesson:FA-C436-arshe}



\begin{quote}
Before the new one: say the Persian for a pier, then the Persian for a captain.
\end{quote}

\subsection*{You'll want to know: \fa{عرشه}}
\textbf{\fa{عرشه}} (\emph{'arshe}) — \textquotedblleft{}a deck (of a ship)\textquotedblright{}.

\begin{quote}
\fa{مسافران} \fa{روی} \fa{عرشه} \fa{ایستاده} \fa{بودند}.
\end{quote}

\begin{quote}
\emph{Mosâferân ru-ye 'arshe istâde budand.}
\end{quote}

\begin{quote}
The passengers were standing on the deck.
\end{quote}

\emph{Ru-ye 'arshe} — on deck.

\begin{cousinweb}
\emph{'arshe} was borrowed from Arabic and is usually traced to Arabic \emph{ʿarš} 'a roof, a raised platform, a throne'. EtymDB does not list it.
\end{cousinweb}

\begin{grammarlens}[title={What you've built}]
One more word for telling, explaining and describing at B1.
\end{grammarlens}

\subsection*{Your turn}
\begin{itemize}
\raggedright
  \item \emph{Say it:} \emph{'arshe}
  \item \emph{Say it:} \emph{'arshe}, once more
  \item \emph{Say it:} a sentence of your own with \emph{'arshe}
  \item \emph{Recall:} say the Persian for a pier, then the Persian for a captain, then say \emph{'arshe} again
\end{itemize}

\begin{tcolorbox}[breakable,colback=teal!4,colframe=teal!35!black,title={Before you move on}]
What does \fa{عرشه} mean? (\textquotedblleft{}A deck (of a ship)\textquotedblright{}.) Say it once more.
\end{tcolorbox}
\section[\texorpdfstring{\fa{لنج} (lenj)}{لنج (lenj)}]{\fa{لنج} (lenj) — a dhow, a wooden cargo boat (of the Gulf)}

\label{lesson:FA-C436-lenj}



\begin{quote}
Before the new one: say the Persian for a captain, then the Persian for a deck.
\end{quote}

\subsection*{You'll want to know: \fa{لنج}}
\textbf{\fa{لنج}} (\emph{lenj}) — \textquotedblleft{}a dhow, a wooden cargo boat (of the Gulf)\textquotedblright{}.

\begin{quote}
\fa{این} \fa{لنج} \fa{از} \fa{بندر} \fa{به} \fa{جزیره} \fa{میرود}.
\end{quote}

\begin{quote}
\emph{In lenj az bandar be jazire miravad.}
\end{quote}

\begin{quote}
This dhow goes from the port to the island.
\end{quote}

Built of wood in the ports of the south, and still in use for trade.

\begin{cousinweb}
The origin of \emph{lenj} is not settled, and EtymDB does not list it. It names the large wooden boats of the Persian Gulf.
\end{cousinweb}

\begin{grammarlens}[title={What you've built}]
One more word for telling, explaining and describing at B1.
\end{grammarlens}

\subsection*{Your turn}
\begin{itemize}
\raggedright
  \item \emph{Say it:} \emph{lenj}
  \item \emph{Say it:} \emph{lenj}, once more
  \item \emph{Say it:} a sentence of your own with \emph{lenj}
  \item \emph{Recall:} say the Persian for a captain, then the Persian for a deck, then say \emph{lenj} again
\end{itemize}

\begin{tcolorbox}[breakable,colback=teal!4,colframe=teal!35!black,title={Before you move on}]
What does \fa{لنج} mean? (\textquotedblleft{}A dhow, a wooden cargo boat (of the Gulf)\textquotedblright{}.) Say it once more.
\end{tcolorbox}
\section[\texorpdfstring{\fa{لنگر} \fa{انداختن} (langar andâkhtan)}{لنگر انداختن (langar andâkhtan)}]{\fa{لنگر} \fa{انداختن} (langar andâkhtan) — to drop anchor, to anchor}

\label{lesson:FA-C436-langar-andakhtan}



\begin{quote}
Before the new one: say the Persian for a deck, then the Persian for a dhow.
\end{quote}

\subsection*{You'll want to know: \fa{لنگر} \fa{انداختن}}
\textbf{\fa{لنگر} \fa{انداختن}} (\emph{langar andâkhtan}) — \textquotedblleft{}to drop anchor, to anchor\textquotedblright{}.

\begin{quote}
\fa{کشتی} \fa{نزدیک} \fa{ساحل} \fa{لنگر} \fa{انداخت}.
\end{quote}

\begin{quote}
\emph{Kashti nazdik-e sâhel langar andâkht.}
\end{quote}

\begin{quote}
The ship dropped anchor near the shore.
\end{quote}

\emph{langar} 'anchor' + \emph{andâkhtan} 'to throw'. Present stem \emph{andâz-}: \emph{langar miandâzad}.

\begin{cousinweb}
\emph{langar andâkhtan} is built from \emph{langar} 'anchor', which EtymDB links with Greek \emph{ánkyra} 'anchor', and \emph{andâkhtan} 'to throw', usually traced to Middle Persian \emph{handāxtan}.
\end{cousinweb}

\begin{grammarlens}[title={What you've built}]
One more word for telling, explaining and describing at B1.
\end{grammarlens}

\subsection*{Your turn}
\begin{itemize}
\raggedright
  \item \emph{Say it:} \emph{langar andâkhtan}
  \item \emph{Say it:} \emph{langar andâkhtan}, once more
  \item \emph{Say it:} a sentence of your own with \emph{langar andâkhtan}
  \item \emph{Recall:} say the Persian for a deck, then the Persian for a dhow, then say \emph{langar andâkhtan} again
\end{itemize}

\begin{tcolorbox}[breakable,colback=teal!4,colframe=teal!35!black,title={Before you move on}]
What does \fa{لنگر} \fa{انداختن} mean? (\textquotedblleft{}To drop anchor, to anchor\textquotedblright{}.) Say it once more.
\end{tcolorbox}
